\subsection{Meager equivalence relations in Polish spaces}

Lemma 3. --- Let \((\mathrm{U}_{i})_{i\in\mathrm{I}}\) be a finite family of nonempty open sets of a topological space X and let O be an open dense subset of \(X\times X\) . There exists a family \((\mathrm{V}_{i})_{i\in\mathrm{I}}\) of nonempty open subsets of X such that \(V_{i}\subset U_{i}\) for all \(i\in I\) and such that \(V_{i}\times V_{i'}\subset O\) for every pair \((i,i')\) of distinct elements of I.

For every pair \((j_{1}, j_{2})\) of distinct elements of I, the set of families \((x_{i}) \in \mathrm{X}^{\mathrm{I}}\) such that \((x_{j_{1}}, x_{j_{2}}) \in \mathrm{O}\) is a dense open subset of \(X^{I}\) . The intersection \(\Omega\) of these open sets, as \((j_{1}, j_{2})\) ranges over the finite set of pairs of distinct elements of I, is therefore a dense open set of \(X^{I}\) .

Therefore, \(\Omega \cap \prod_{i\in \mathbf{I}}\mathrm{U}_i\) is a nonempty open set of \(\mathrm{X}^{\mathrm{I}}\) , hence contains an open set of \(\mathrm{X}^{\mathrm{I}}\) of the form \(\prod_{i\in \mathbf{I}}\mathrm{V}_i\) , where for every \(i\in \mathbf{I}\) , \(\mathrm{V}_i\) is a nonempty open subset of \(\mathrm{U}_i\) . This proves the lemma.

PROPOSITION 8. --- Let P be a nonempty Polish topological space and let R be an equivalence relation on P whose graph is a meager subset of P × P. There exists a continuous injective map from the set \(\{0,1\}^{\mathbf{N}}\) , endowed with the product topology of the discrete topologies, into P whose image meets each equivalence class under R in at most one point.

Endow the space P with a metric \(d\) compatible with its topology for which it is complete. Let \((\mathrm{A}_n)_{n\in \mathbf{N}}\) be an increasing sequence of closed subsets of \(\mathrm{P}\times \mathrm{P}\) , with empty interiors, such that the graph \(\Gamma_{\mathrm{R}}\) of the relation R is contained in the union of the \(\mathrm{A}_n\) .

Let \(\mathcal{O}\) be the set of nonempty open subsets of P. We shall construct by induction a sequence \((f_n)_{n \in \mathbf{N}}\) , where for every \(n\) , \(f_n\) is a map from \(\{0,1\}^n\) into \(\mathcal{O}\) , satisfying the following properties:

\begin{enumerate}
\def\labelenumi{(\roman{enumi})}
\item
  For every \(n \geqslant 1\) , every \(x \in \{0,1\}^{n-1}\) and every \(t \in \{0,1\}\) , the closure of the open set \(f_n(x,t)\) is contained in \(f_{n-1}(x)\) ;
\item
  For every \(x \in \{0,1\}^n\) , we have \(\mathrm{diam}(f_n(x)) \leqslant 2^{-n}\) ;
\item
  For every \(n \geqslant 1\) and every pair \((x, x')\) of distinct elements of \(\{0, 1\}^n\) , \(f_n(x) \times f_n(x')\) does not meet \(\mathrm{A}_{n-1}\) .
\end{enumerate}

Choose a nonempty open set U of P of diameter \(\leqslant 1\) and one defines \(f_{0}\) as the constant map with image \(\{U\}\) . Suppose the maps \(f_{0}, \ldots, f_{n}\) have been constructed.

Let \(p \colon \{0,1\}^{n+1} \to \{0,1\}^n\) be the map defined by \(p(x_0, \ldots, x_n) = (x_0, \ldots, x_{n-1})\) . By Lemma 3 above, applied to the family of open sets \((f_n(p(x)))\) for \(x \in \{0,1\}^{n+1}\) and to the dense open set \(\mathbb{C}\mathrm{A}_n\) of \(\mathrm{P} \times \mathrm{P}\) , there exists a family \((g(x))_{x \in \{0,1\}^{n+1}}\) of nonempty open subsets of P such that \(g(x) \subset f_n(p(x))\) for all \(x\) and such that \((g(x) \times g(x')) \cap \mathrm{A}_n\) is empty for every pair \((x,x')\) of distinct elements of \(\{0,1\}^{n+1}\) . One then defines the map \(f_{n+1}\) by choosing, for each element \(x \in \{0,1\}^{n+1}\) , a nonempty open subset of \(g(x)\) whose diameter is \(\leqslant 2^{-n-1}\) and whose closure is contained in \(g(x)\) .

For every element \(x = (x_n)_{n \in \mathbf{N}}\) of \(\{0,1\}^{\mathbf{N}}\) , the sequence of sets \((f_n(x_0, \ldots, x_{n-1}))_{n \in \mathbf{N}}\) is a decreasing sequence of open subsets of P each of which contains the closure of the next and whose diameter tends to 0; the intersection of this sequence of sets is therefore reduced to a point (TG, II, p. 15), which one denotes by \(f(x)\) . If two points \(x\) , \(x'\) of \(\{0,1\}^{\mathbf{N}}\) satisfy \(x_i = x'_i\) for \(i \leqslant n\) , we have \(d(f(x), f(x')) \leqslant 2^{-n}\) . Consequently, the map \(f: \{0,1\}^{\mathbf{N}} \to \mathbf{P}\) is continuous. Let \(x\) and \(x'\) be distinct elements of \(\{0,1\}^{\mathbf{N}}\) . For \(n \in \mathbf{N}\) such that \((x_0, \ldots, x_n) \neq (x_0', \ldots, x_n')\) , the open set \(f_{n+1}(x_0, \ldots, x_n) \times f_{n+1}(x_0', \ldots, x_n')\) is disjoint from \(A_n\) , by definition of \(f_{n+1}\) , so the pair \((f(x), f(x'))\) does not belong to \(A_n\) . It follows that \(f(x)\) and \(f(x')\) are not equivalent for the relation R. Consequently, \(f\) is injective and its image meets each equivalence class under R in at most one point. The proposition is thus proved.

